\documentclass[conference]{IEEEtran}
\IEEEoverridecommandlockouts

\usepackage{cite}
\usepackage{amsmath,amssymb,amsfonts}
\usepackage{algorithmic}
\usepackage{graphicx}
\usepackage{textcomp}
\usepackage{xcolor}
\usepackage{booktabs}
\usepackage{url}

\begin{document}

\title{Adaptive Evacuation with Graph-Based Deep Reinforcement Learning for Dynamic Indoor Hazards}

\author{\IEEEauthorblockN{Author Name(s)}
\IEEEauthorblockA{Affiliation\\
City, Country\\
email / ORCID}}

\maketitle

\begin{abstract}
% TODO: Rewrite only after the controlled experiments are complete.
This abstract will be written from the verified final experimental results. It will state the evacuation problem, the identified research gap, the proposed graph-based Double DQN framework, the experimental comparison, the principal findings, and the observed limitations without unsupported generalization claims.
\end{abstract}

\begin{IEEEkeywords}
Indoor Evacuation, Graph Neural Networks, Deep Reinforcement Learning, Double DQN, Dynamic Hazards, Intelligent Environments
\end{IEEEkeywords}

\section{Introduction}
\label{sec:introduction}

% TODO: Build from the verified research problem and literature review.
% Do not reuse unsupported claims from legacy_draft_v0.

\section{Related Work}
\label{sec:related_work}

% TODO: Review classical path planning, DRL evacuation, CNN-based navigation,
% graph-based navigation, and recent GNN-RL evacuation studies.

\section{Problem Formulation and Proposed Method}
\label{sec:methodology}

% TODO: Write from the audited implementation.
% Required subsections: MDP, environment, grid-to-graph, node features,
% hazard-aware reward/reference planning, GCN encoder, agent-centric embedding,
% Double-DQN target, and training procedure.

\section{Experimental Methodology}
\label{sec:experiments}

% TODO: Define benchmark layouts, training/evaluation splits, baselines,
% seeds, metrics, and controlled experiments before inserting results.

\section{Results and Analysis}
\label{sec:results}

% TODO: Insert only reproducible experimental results.
% Planned analysis: baseline comparison, generalization, scaling, hazard
% exposure, path efficiency, ablations, and failure cases.

\section{Discussion and Conclusion}
\label{sec:conclusion}

% TODO: Discuss strengths, limitations, scaling behavior, implications,
% and future work. Do not claim safety guarantees or universal deployment.


\bibliographystyle{IEEEtran}
\bibliography{references}

\end{document}
